
\section{Movimento até uma posição alvo}
\label{sec:imp_move}

O movimento do manipulador até uma posição pretendida envolve a transformação das coordenadas do alvo, o cálculo da configuração articular e a execução do movimento através das funções desenvolvidas em MATLAB.

Quando a posição do alvo é fornecida no referencial do mundo do CoppeliaSim, o procedimento pode ser dividido nas seguintes etapas:

\begin{enumerate}
    \item \textbf{Transformação de coordenadas:}
    a posição pretendida $(X_{\mathrm{target}},Y_{\mathrm{target}},Z_{\mathrm{target}})$ é convertida para o referencial da base do manipulador através da função \texttt{aux.world\_to\_base}, considerando a translação e a rotação entre os dois referenciais;

    \item \textbf{Cálculo da configuração articular:}
    a função \texttt{inverse\_kin} determina os ângulos necessários para alcançar a posição pretendida, respeitando as restrições geométricas e articulares definidas. Caso não exista uma solução válida, o movimento pode ser rejeitado antes da sua execução;

    \item \textbf{Execução do movimento:}
    após determinar uma configuração válida, o manipulador pode ser deslocado através da função \texttt{move\_joints}, que realiza a interpolação dos ângulos articulares, ou da função \texttt{move\_linear}, utilizada quando se pretende seguir uma trajetória aproximadamente retilínea no espaço cartesiano;

    \item \textbf{Verificação do posicionamento:}
    os ângulos articulares obtidos durante a execução podem ser utilizados pela função \texttt{direct\_kin} para calcular a posição do centro de preensão e avaliar a diferença relativamente à posição pretendida.
\end{enumerate}

\subsection{Movimento no espaço das juntas}

A função \texttt{move\_joints} permite deslocar o manipulador entre duas configurações articulares, garantindo uma evolução suave dos ângulos das juntas.

Em vez de enviar diretamente os ângulos finais, a função calcula, em cada passo de simulação, uma configuração intermédia entre a configuração inicial e a pretendida. Todas as juntas iniciam e terminam o movimento ao mesmo tempo, independentemente da amplitude da rotação de cada uma.

A evolução entre as duas configurações segue um perfil suave, em forma de S: as juntas arrancam devagar, aceleram até ao meio do movimento e abrandam até parar na posição pretendida. Desta forma, o movimento começa e termina com velocidade nula, evitando alterações bruscas da velocidade e solavancos no manipulador.

A duração do movimento é definida pela junta que tem de rodar mais. Esta duração é calculada de modo a que a velocidade dessa junta nunca ultrapasse a velocidade máxima pretendida, de $30^\circ$/s por defeito. As restantes juntas, como rodam menos no mesmo tempo, movem-se mais devagar.

A execução é realizada de forma sincronizada, sendo enviadas configurações articulares intermédias ao manipulador em cada passo de simulação.

No final, a função verifica se todas as juntas se encontram a menos de $1^\circ$ dos respetivos valores pretendidos.
\subsection{Movimentos em linha reta}

Para movimentos de aproximação a objetos, nomeadamente durante a descida da garra até uma peça, é utilizada a função \texttt{move\_linear}.

Ao contrário do movimento no espaço das juntas, esta função faz evoluir um ponto de referência ao longo da direção que une a posição atual do centro de preensão à posição pretendida.

Em cada passo de simulação, o programa atualiza a velocidade do ponto de referência, respeitando a velocidade máxima definida e a aceleração configurada através do parâmetro \texttt{cfg.ACC}.

A velocidade é aumentada gradualmente e reduzida à medida que o ponto se aproxima do destino, de forma a evitar uma paragem brusca.

Para cada posição intermédia, a função \texttt{inverse\_kin} calcula os ângulos articulares correspondentes, que são enviados ao manipulador.

Após o ponto de referência atingir o destino, o programa compara a posição pretendida com a posição calculada pela cinemática direta a partir dos ângulos efetivamente lidos.

Quando necessário, é aplicada uma correção progressiva ao ponto de referência para compensar diferenças de posicionamento, incluindo os efeitos associados à gravidade.

O movimento é considerado concluído com sucesso quando o erro de posicionamento é inferior a $3$ mm, dentro das condições de execução estabelecidas.

\subsection{Planeamento de movimentos verticais}

Antes de executar determinados movimentos verticais, a função \texttt{plan\_vertical} seleciona uma inclinação fixa da garra que permita percorrer toda a trajetória.

Este procedimento é particularmente importante durante a aproximação a objetos situados perto da base do manipulador, onde pequenas alterações de posição podem exigir variações significativas dos ângulos articulares.

O algoritmo testa diferentes inclinações, começando pela orientação vertical para baixo, correspondente a $\psi=90^\circ$.

Para cada inclinação candidata, são verificadas sucessivamente as posições ao longo da trajetória vertical, em intervalos de $2$ mm.

Uma inclinação apenas é aceite quando são satisfeitas as seguintes condições:

\begin{itemize}
    \item Existe uma solução válida de cinemática inversa para todos os pontos da trajetória;

    \item A variação de cada uma das juntas J1 a J4 entre posições consecutivas não ultrapassa $5^\circ$;

    \item Quando a verificação de colisões está ativada, as configurações articulares não provocam colisões do manipulador consigo próprio.
\end{itemize}

Durante esta verificação, é utilizada a configuração com o cotovelo para cima, sendo mantida a mesma inclinação da garra ao longo de todo o movimento.

Esta abordagem reduz a possibilidade de alterações abruptas das configurações articulares durante a descida.

A Figura~\ref{fig:move} apresenta uma representação esquemática da sequência de aproximação a uma peça.

\begin{figure}[htbp]
    \centering
    % Sequência do movimento até uma peça, vista lateral (esboço, sem escala)
\begin{tikzpicture}[scale=1.2, >=stealth,
    elo/.style={line width=2pt, line cap=round},
    junta/.style={circle, draw, fill=white, line width=0.9pt, inner sep=1.6pt},
    pt/.style={circle, fill, inner sep=1.5pt},
    every node/.style={font=\footnotesize}]

  % mesa e base
  \fill[pattern=north east lines, pattern color=gray!60] (-1.2,-0.25) rectangle (10.2,0);
  \draw[thick] (-1.2,0) -- (10.2,0) node[below left=2pt and 0pt] {mesa};
  \fill[gray!30] (-0.45,0) rectangle (0.45,0.5);
  \draw (-0.45,0) rectangle (0.45,0.5);
  % braço na pose inicial (esboço)
  \coordinate (J2) at (0,1.6);
  \draw[elo, gray!70] (0,0.5) -- (J2) -- (0.3,3.3) -- (1.6,4.1);
  \node[junta] at (J2) {}; \node[junta] at (0.3,3.3) {};
  % peça
  \fill[red!65] (6.3,0) rectangle (7.1,0.8);
  \draw (6.3,0) rectangle (7.1,0.8);

  \coordinate (P0) at (1.6,4.1);
  \coordinate (P1) at (6.7,3.0);
  \coordinate (P2) at (6.7,1.4);

  % 1: move_joints (trajetória curva da garra)
  \draw[->, blue!70!black, line width=1.1pt, dashed] (P0) .. controls (3.6,5.2) and (5.8,4.4) .. (P1);
  \node[blue!70!black, align=center] at (3.9,5.35) {\textbf{1.} \texttt{move\_joints}};
  % 2: move_linear (descida vertical)
  \draw[->, orange!85!black, line width=1.3pt] (P1) -- (P2);
  \foreach \h in {2.8,2.6,...,1.5} \draw[orange!85!black] (6.6,\h) -- (6.8,\h);
  \node[orange!85!black, align=left, anchor=west] at (7.7,2.4) {\textbf{2.} \texttt{move\_linear}\\descida em linha reta,\\uma só inclinação\\(\texttt{plan\_vertical},\\de 2 em 2 mm)};
  % 3: subida e go_home
  \draw[->, green!45!black, line width=1.1pt, dotted] (6.2,1.5) -- (6.2,2.9) .. controls (5.0,4.0) and (3.0,4.6) .. (1.8,4.25);
  \node[green!45!black, align=center] at (3.5,3.1) {\textbf{3.} \texttt{go\_home}: sobe na\\vertical e regressa};

  \node[pt] at (P0) {}; \node[above left=1pt and 0pt] at (P0) {pose inicial};
  \node[pt] at (P1) {}; \node[above right=5pt and 2pt] at (P1) {aproximação};
  \node[pt] at (P2) {};
  \draw[<->, gray!70!black, thin] (7.35,0) -- (7.35,1.4) node[midway, right, black] {2 cm};
  \draw[gray!70!black, thin, dashed] (6.7,1.4) -- (7.45,1.4);
\end{tikzpicture}
    \caption{Representação esquemática da sequência de movimento até uma peça, numa vista lateral e sem escala.}
    \label{fig:move}
\end{figure}

\subsection{Regresso à pose inicial}

A função \texttt{go\_home} permite deslocar o manipulador para uma configuração inicial previamente definida, independentemente da posição articular de partida, desde que o movimento seja realizável.

O procedimento começa pela abertura da garra e pela leitura dos ângulos articulares atuais.

Através da função \texttt{direct\_kin}, é determinada a posição do centro de preensão. Caso a sua altura seja inferior ao valor de segurança definido, de $0{,}15$ m relativamente à base, o algoritmo procura realizar uma subida gradual antes do regresso à configuração inicial.

Durante esta subida, são calculadas sucessivamente novas configurações articulares através da função \texttt{inverse\_kin}, procurando manter a inclinação inicial da garra.

Se não for possível encontrar uma configuração válida ou se forem detetadas variações articulares excessivas, a subida é interrompida e o programa avança para a etapa seguinte.

Posteriormente, a função \texttt{move\_joints} é utilizada para executar o movimento suave até à configuração inicial.

Por fim, os ângulos articulares são novamente lidos, sendo o regresso considerado bem-sucedido quando todas as juntas apresentam um desvio inferior a $1^\circ$ relativamente aos valores pretendidos.

Este procedimento baseia-se nas leituras articulares, permitindo utilizar a mesma lógica de controlo com a interface do manipulador simulado ou com uma interface compatível de um braço físico.
